\section{Conclusion and Limitations}
\label{sec:conclusion}

We formulated a sparse-to-tube view of training-free video copy localization and a way to
test it. The pipeline proposes sparse correspondences, fits locally linear temporal tracks with
support on both axes, and restricts high-resolution alignment to finite tubes. The formulation
makes the cost structure explicit: fixed-stride coarse matching remains $\Theta(MN)$, a finite
tube gives $O(LW)$ fine alignment per path under direct offset addressing, and an implementation
that scans a predecessor row turns the same computation into $O(PLW^2)$. That difference is not
cosmetic; it separates an arithmetic argument from a wall-clock result.

The paper is organized around three hypotheses rather than an accuracy claim. Candidate
coverage (H1) can fail even with perfect refinement, because an uncovered copy lies outside
every tube. Boundary recovery (H2) can hold while the system becomes slower, because a smaller
domain is not automatically a smaller measured cost (H3). Making these failures explicit tells
a reader in advance which observation would refute the design and which comparison would expose
a hidden cost. The novelty gate is deliberately severe: a baseline of one grid, mutual Top-$K$,
RANSAC and standard Smith--Waterman already contains every classical ingredient, and if the
additional mechanisms do not beat it at equal dot-product budget, the correct report is a
negative result.

\paragraph{Limitations.}
Sparse sampling cannot propose a copy shorter than $(h_{\min}-1)s+1$ feature frames with a
stride-$s$ grid, and short copies are where deduplication systems are most sensitive. The slope
range is positive and bounded, so reversed or extreme speed changes are outside the model;
non-linear speed changes are only approximated by multiple local lines. Repeated backgrounds
and static logos can generate many mutual anchors at different reference times; the dual-axis
support test suppresses tracks pinned to one instant but cannot remove all pseudo-evidence. The
number of tubes can grow with spurious tracks, offsetting the alignment saving. The score $q$
is uncalibrated and not comparable to detector objectness, so thresholds must be per method.
Finally, VCSL provides interval boxes, not dense correspondence, so neither the official metric
nor our diagnostics establish frame-accurate correspondence.

\paragraph{What remains.}
No benchmark result is reported. A zero-error complete validation run, a frozen threshold for
every method, a single test execution with four denominators, a controlled comparison against
the Simple Sparse-to-Band baseline on identical candidates, a real-hardware CUDA parity and
timing study, and a qualitative figure on real pairs are outstanding. Until then the honest
status is a specified method and a falsifiable protocol, not a demonstrated improvement. If the
frozen runs support the mechanism, the contribution is a quality--cost trade-off under fixed
descriptors; if they do not, the same protocol will have produced a useful negative result about
where sparse proposals stop being sufficient.
